\documentclass[10pt,a4paper]{amsart}
\usepackage[margin=19mm]{geometry}
\usepackage{amsmath,amssymb,mathtools}
\usepackage[scr=rsfs,cal=euler]{mathalfa}
\usepackage{xfrac}
\usepackage[unicode,hidelinks,bookmarks=false]{hyperref}
\newcommand{\ie}{i.e.\ }
\newcommand{\cf}{cf.\ }
\newcommand{\resp}{respectively\ }
\newcommand{\Subsection}{\subsection}
\providecommand{\nobreakdash}{\nobreak\hbox{-}}
\setlength{\emergencystretch}{2em}
\multlinegap0pt
\usepackage{tikz}
\usepackage{amssymb}
\usepackage{enumerate}
\usepackage{float}
\usepackage{bm}
\usepackage{algorithm}\usepackage{algpseudocode}
\usepackage{mathtools}
\newtheorem{corollary}{Corollary}
\newtheorem{theorem}{Theorem}
\usepackage{cleveref}
\crefname{corollary}{Corollary}{Corollarys}
\crefname{theorem}{Theorem}{Theorems}
\newcommand{\A}{\mathcal{A}}
\DeclareMathOperator{\FMset}{\mathcal{F}}
\newcommand{\N}{\mathbb{N}}
\newcommand{\R}{\mathbb{R}}
\newcommand{\T}{\intercal}
\newcommand{\esssupp}{\mathrm{esssupp}}
\newcommand{\mb}[1]{\mathbb #1}
\DeclareMathOperator{\relint}{relint}
\newcommand{\sample}{ X_1,\dots, X_n}
\renewcommand{\vec}[1]{\mathbf{#1}}
\title{On the Uniqueness of Fr\'echet Means for Polytope Norms}
\author{Roan Talbut, Andrew McCormack, Anthea Monod}
\date{Source: arXiv:2603.15785v1 (CC BY 4.0)}
\begin{document}
\maketitle

\noindent\emph{Source and licence.} On the Uniqueness of Fr\'echet Means for Polytope Norms. Version arXiv:2603.15785v1 (CC BY 4.0), distributed by arXiv under the Creative Commons Attribution 4.0 International licence, http://creativecommons.org/licenses/by/4.0/, as recorded in the arXiv licence metadata for this exact version and shown on the abstract page https://arxiv.org/abs/2603.15785v1. Full licence notice: \texttt{licenses/arxiv-2603.15785-CC-BY-4.0.txt}.\par\smallskip

\noindent\textit{Editorial scope.} Counted unit: the article's theorem thm:uniqueness\_prob\_convergence (source0.tex lines 678--680). The article's complete original proof of it is delivered with it (source0.tex lines 681--718, one \texttt{proof} environment of 38 lines). No mathematics is written, repaired or re-derived: the statement and the proof are the article's own lines, in the article's own notation, and no retained line is substituted or re-flowed.  Retained uncounted as the source-local context the proof needs: lines 564--568; lines 668--670.  Nothing was omitted from any retained span, so the package carries no omission record.  No retained line contains a citation, so the wrapper carries no bibliography and no bibliography entry.  The retained lines contain the article's own diagram code, which the wrapper typesets with the standard tikz packages; no image file is delivered and the package declares no asset dependency.  Editorial formatting only, in addition: an AMS \texttt{amsart} preamble that restates the article's own declarations and macros used by the retained lines, namely the environments \texttt{corollary}, \texttt{theorem} on separate counters, together with the standard packages those lines need.  The retained environments print in this document as Corollary 1, Theorem 1 in order of appearance, whereas the article prints the same items with the numbering of its own full document.  The complete native source of the article as distributed in the arXiv e-print for this exact version is bundled unchanged as \texttt{sources/196499/source0.tex}.
	\begin{corollary}
		\label{cor:FMisPolytope}
		Fix a normed space $(\R^k, \| \cdot \|_B)$ with unit ball $B$. There exist unique $d_i \in \R_{\geq 0}$ such that $$\FMset (\vec x_1,\ldots,\vec x_n) = \bigcap_{i \in [n]}\{d_iB + \vec x_i\} = \bigcap_{i \in [n]}\{d_i\partial B + \vec x_i\},
		$$ where $\partial B$ is the boundary of $B$. In particular, if $B$ is a polytope then $\FMset(\vec x_1, \dots, \vec x_n)$ is a polytope. Moreover, for every $i \in [n]$ there exists a proper face $F_i$ of $B$ such that $\FMset$ is given by $\bigcap_{i \in [n]}\{d_iF_i + \vec x_i\}$. 
	\end{corollary}

	\begin{theorem}\label{thm:consistent_FM}
		Let $\{X_i\}_{i \in \N}$ be a sequence of i.i.d. random vectors with $X_i \sim \mb{P}$, where $\mb P$ has a unique population Fr\'echet mean $\theta_*$. Let  $\theta_n$ be a random variable such that $\theta_n \in \FMset(X_1, \dots, X_n)$ almost surely. Then $\theta_n$ converges almost surely to $\theta_*$.
	\end{theorem}

	\begin{theorem}\label{thm:uniqueness_prob_convergence}
		Suppose $X \sim \mb P$, with $\esssupp(\mb P) = \R^k$. Let $p_n$ be the probability that $\FMset(X_1,\dots,X_n)$ consists of a single element. Then $p_n$ converges to 1 as $n \rightarrow \infty$.
	\end{theorem}

	\begin{proof}
		Suppose that there is some sample Fr\'echet mean $\theta \in \FMset(\sample)$ such that for every possible active constraint $\vec a_i \in A$, there is some data point $X_{j(i)}$ satisfying $\A(X_{j(i)}-\theta) = \vec a_i$. Then by \Cref{cor:FMisPolytope}, the Fr\'echet mean set is contained in the intersection of corresponding faces, $\bigcap_i \|\theta-X_{j(i)}\| F_i+X_{j(i)}$ and so
		\begin{align}
			\FMset(X_1,\dots,X_n) \subseteq \{ \vec x: \vec a_i^{\T}(\vec x-X_{j(i)}) = \|\theta-X_{j(i)}\| \} \label{eq:FM_uniqueness_linear_constraints}.
		\end{align}
		As the set of all $\vec a_i$ spans $\R^k$, the linear constraints in \ref{eq:FM_uniqueness_linear_constraints} define a unique point in $\FMset(X_1, \dots, X_n)$. It now suffices to show that, with probability tending to 1, there is some Fr\'echet mean $\theta$ such that $A = \{ \A(X_j - \theta), j=1,\dots,n \}$.
		
		For each facet $F_i$, consider the intersection $S_i = \bigcap_{\theta \in B_1(\theta_*)} (\theta - \relint (\R_+ F_i)) \subset \R^k$ (see \Cref{fig:unique_prob_to_1}). This is the set of points $\vec x$ with $\A(\vec x - \theta) = \vec a_i$ for all $\theta$ in the unit ball about $\theta_*$. This intersection is full dimensional, so $\mb P (X_j \in S_i) = c_i > 0$. Let $b_n$ be the probability that there is some $\theta \in \FMset (X_1, \dots, X_n)$ in $B_1(\theta_*)$. By the argument above, we can bound $p_n$ below using $b_n, c_i$:
		\begin{align*}
			p_n &\geq \mb P(A = \{ \A(X_j - \theta), j=1,\dots,n \}) \\
			&\geq \mb P(\{\emptyset \neq \FMset(X_1,\dots,X_n) \cap B_1(\theta_*)\} \cap \{ \forall C_i, \emptyset \neq C_i \cap \{ X_1, \dots, X_n \} \}) \\
			&\geq 1 - (1-b_n) - \textstyle \sum_{i \in [r]} (1-c_i)^n
		\end{align*}
		By the consistency of $\theta_n$ (\Cref{thm:consistent_FM}), $b_n \rightarrow 1$. So this lower bound converges to 1.
		\begin{figure}[ht!]
			\centering
			\begin{tikzpicture}
				\draw (0,-1) -- (1, 0) -- (0,1) -- (-1,0) -- (0,-1);
				\draw (-6,-1) -- (6,-1);
				\draw (-6,1) -- (6,1);
				\draw (-1,-3) -- (-1,3);
				\draw (1,-3) -- (1,3);
				\draw[<->] (-1,0) -- (0,0) node[anchor=north east] {$1$};
				\filldraw (0,0) circle (2pt) node[anchor=north west] {\small$\theta_*$};
				\filldraw (-0.3,0.4) circle (2pt) node[anchor=south east] {\small$\theta_n$};
				\node at (4,2) {$S_2$};
				\node at (-4,2) {$S_1$};
				\node at (4,-2) {$S_3$};
				\node at (-4,-2) {$S_4$};
				\filldraw (2,2) circle (2pt) node[anchor=south west] {\small$X_{j(2)}$};
				\filldraw (-1.8,-2) circle (2pt) node[anchor=north east] {\small$X_{j(4)}$};
				\filldraw (2,-2) circle (2pt) node[anchor=north west] {\small$X_{j(3)}$};
				\filldraw (-2.4,2.3) circle (2pt) node[anchor=south east] {\small$X_{j(1)}$};
			\end{tikzpicture}
			\caption{A visualisation of the cones $S_i$ (defined in the proof of \cref{thm:uniqueness_prob_convergence}) in the case of the $\ell_1$ norm, and a sample such that every cone contains some data point. As each cone contains a data point, whenever there is a Fr\'echet mean within the unit ball of $\theta_*$, it must be a unique Fr\'echet mean.}
			\label{fig:unique_prob_to_1}
		\end{figure}
	\end{proof}

\end{document}
